% !TeX root = ../navier-stokes-manuscript.tex
% ============================================================
% 2.1 Vortex Stretching
% ============================================================

Axial stretching lengthens one narrow material vortex-tube segment; because its material volume stays fixed, its cross-section contracts, its rotating fluid moves inward, and the interior spins faster. The velocity field carrying that segment obeys
\[
  \underbrace{\partial_t u}_{\text{local acceleration}}
  +
  \underbrace{(u\cdot\nabla)u}_{\text{nonlinear transport}}
  =
  \underbrace{-\nabla p}_{\text{pressure}}
  +
  \underbrace{\nu\Delta u}_{\text{viscous diffusion}}.
\]
Taking the curl removes pressure and turns nonlinear transport into vorticity advection and stretching.

\subsection{Vortex Stretching}\label{sub:ThePerfectlyStretchingVortex}

Locally, the segment is axisymmetric, its strain is uniform across the narrow cross-section, and the stretching direction is aligned with its axis. If \(e\) is that axis and \(L(t)\) is the segment's length, then
\[
  S:=\frac12\bigl(\nabla u+(\nabla u)^{\mathsf T}\bigr),
  \qquad
  Se=\sigma(t)e,
  \qquad
  \sigma(t)>0,
  \qquad
  \frac{\dot L(t)}{L(t)}
  =
  e\cdot Se
  =
  \sigma(t).
\]
The positive axial strain rate lengthens the segment. Incompressibility fixes its material volume \(\mathcal V\), so its effective transverse area is \(A(t):=\mathcal V/L(t)\) and its effective radius is \(r_{\mathrm{eff}}(t):=\sqrt{A(t)/\pi}\). Thus
\[
  \nabla\cdot u=0,
  \qquad
  \frac{d}{dt}(A L)=0,
  \qquad
  \frac{\dot A}{A}=-\sigma(t),
  \qquad
  \frac{\dot r_{\mathrm{eff}}}{r_{\mathrm{eff}}}
  =
  -\frac{\sigma(t)}{2}.
\]
The logarithmic derivatives of area and radius are negative, while their contraction magnitudes are \(\sigma(t)\) and \(\sigma(t)/2\). With \(\omega:=\nabla\times u\), the vorticity carried by the segment satisfies
\[
  \frac{D\omega}{Dt}
  =
  S\omega+\nu\Delta\omega,
  \qquad
  \frac{D}{Dt}
  :=
  \partial_t+u\cdot\nabla,
\]
and under the alignment \(\omega=|\omega|e\),
\[
  S\omega
  =
  \sigma(t)\omega,
  \qquad
  \frac12\frac{D|\omega|^2}{Dt}
  =
  \sigma(t)|\omega|^2
  +
  \nu\,\omega\cdot\Delta\omega.
\]
Aligned axial stretching contributes the positive quantity \(\sigma(t)|\omega|^2\), while the viscous contribution has the sign of \(\omega\cdot\Delta\omega\). The squared vorticity grows exactly where their signed sum is positive. The energy identity proved in \Cref{prop:ClassicalKineticEnergyIdentity} and the antisymmetric part of \(\nabla u\) give
\[
  \norm{u(t)}_2
  \leq
  \norm{u_0}_2
  <\infty,
  \qquad
  \left|\frac12\bigl(\nabla u-(\nabla u)^{\mathsf T}\bigr)\right|_{\mathrm F}
  =
  \frac{|\omega|}{\sqrt2}
  \leq
  |\nabla u|_{\mathrm F}.
\]
Finite kinetic energy can coexist with concentration of the rotational velocity gradient inside the narrowing segment. That limit of global control creates the need for a full-field comparison at finer width.

\divider{\NavierStokes{} Scaling}

At a time \(t_0<T_*\), let \(\ell:=r_{\mathrm{eff}}(t_0)\). For \(\lambda>1\), another complete \NavierStokes{} solution at the finer width \(\lambda^{-1}\ell\) is
\[
  u_\lambda(x,t)
  :=
  \lambda u(\lambda x,\lambda^2t),
  \qquad
  p_\lambda(x,t)
  :=
  \lambda^2p(\lambda x,\lambda^2t).
\]
At time \(\lambda^{-2}t_0\), the corresponding segment in \(u_\lambda\) has effective radius \(\lambda^{-1}\ell\). This rescaled field is another solution, not a later material state of the original segment. Its momentum terms transform as
\[
\begin{aligned}
  \partial_tu_\lambda(x,t)
  &=
  \lambda^3(\partial_tu)(\lambda x,\lambda^2t),
  \\
  (u_\lambda\cdot\nabla)u_\lambda(x,t)
  &=
  \lambda^3\bigl((u\cdot\nabla)u\bigr)(\lambda x,\lambda^2t),
  \\
  \nabla p_\lambda(x,t)
  &=
  \lambda^3(\nabla p)(\lambda x,\lambda^2t),
  \\
  \nu\Delta u_\lambda(x,t)
  &=
  \lambda^3\nu(\Delta u)(\lambda x,\lambda^2t).
\end{aligned}
\]
Every momentum contribution receives the same cubic factor, so finer scale gives viscosity no relative advantage over nonlinear transport.

\divider{Critical Height}

Homogeneous Sobolev norms respond differently according to derivative order. Let \(\Lambda:=(-\Delta)^{1/2}\). Since
\[
  \widehat{u_\lambda}(\xi,t)
  =
  \lambda^{-2}\widehat u(\lambda^{-1}\xi,\lambda^2t),
\]
a change of variables gives
\[
  \norm{u_\lambda(t)}_{\dot H^s}^2
  =
  \lambda^{2s-1}
  \norm{u(\lambda^2t)}_{\dot H^s}^2.
\]
Kinetic energy loses one power of \(\lambda\) at \(s=0\), while first derivatives gain one at \(s=1\). The exponent vanishes at \(s=\tfrac12\), defining the scale-neutral critical height
\[
  H(t)
  :=
  \frac12\norm{u(t)}_{\dot H^{1/2}}^2
  =
  \frac12\norm{\Lambda^{1/2}u(t)}_2^2.
\]
Writing \(H_\lambda\) for the same quantity evaluated on \(u_\lambda\),
\[
  \norm{u_\lambda(t)}_2^2
  =
  \lambda^{-1}\norm{u(\lambda^2t)}_2^2,
  \qquad
  H_\lambda(t)=H(\lambda^2t),
  \qquad
  \norm{\nabla u_\lambda(t)}_2^2
  =
  \lambda\norm{\nabla u(\lambda^2t)}_2^2.
\]
Scale neutrality fixes the height under rescaling, but it does not fix whether that height rises or falls in time.

\divider{Nonlinear Production versus Viscous Dissipation}

At critical height, nonlinear transport contributes the signed production
\[
  \mathcal P_{1/2}[u](t)
  :=
  -\operatorname{Re}
  \pair
    {\Lambda^{1/2}u(t)}
    {\Lambda^{1/2}\bigl((u(t)\cdot\nabla)u(t)\bigr)}_{L^2}.
\]
Incompressibility removes the pressure pairing, while viscosity removes height at the nonnegative rate \(\nu\norm{\Lambda^{3/2}u}_2^2\). Consequently,
\[
  H'(t)
  +
  \nu\norm{\Lambda^{3/2}u(t)}_2^2
  =
  \mathcal P_{1/2}[u](t).
\]
The height rises when nonlinear production exceeds viscous dissipation and falls when dissipation is larger. Under rescaling,
\[
  \mathcal P_{1/2}[u_\lambda](t)
  =
  \lambda^2\mathcal P_{1/2}[u](\lambda^2t),
  \qquad
  \nu\norm{\Lambda^{3/2}u_\lambda(t)}_2^2
  =
  \lambda^2\nu\norm{\Lambda^{3/2}u(\lambda^2t)}_2^2.
\]
Both acquire the same quadratic multiplier, leaving the duration of a nonlinear transition unresolved.

\divider{The Heat Clock}

The heat equation \(v_t=\nu\Delta v\) supplies a linear comparison for that duration:
\[
  \widehat v(\xi,t+\delta t)
  =
  e^{-\nu|\xi|^2\delta t}\widehat v(\xi,t).
\]
A scale-localized velocity structure of width \(r\) occupies frequencies \(|\xi|\simeq r^{-1}\) and changes by order one on the comparison time determined by
\[
  \nu|\xi|^2\delta t\simeq1,
\]
so the comparison time is
\[
  \tau_\nu(r)
  :=
  \frac{r^2}{\nu}.
\]
This linear heat time does not realize the nonlinear transition. At a finer width,
\[
  \tau_\nu(\lambda^{-1}r)
  =
  \lambda^{-2}\tau_\nu(r).
\]
The comparison time contracts by the same power as time under the full \NavierStokes{} rescaling.

\divider{The Zeno Cascade}

For the dyadic widths and their heat times
\[
  r_n:=2^{-n}r_0,
  \qquad
  \tau_n:=\frac{r_n^2}{\nu},
\]
the times shrink geometrically:
\[
  \tau_n
  =
  \frac{r_0^2}{\nu}4^{-n},
  \qquad
  \sum_{n=0}^{\infty}\tau_n
  =
  \frac{4r_0^2}{3\nu}
  <
  \infty.
\]
The finite geometric sum supplies only a possible schedule. Its realization requires one \NavierStokes{} velocity field to exhibit scale-localized velocity structures at widths
\[
  r_0>r_1>r_2>\cdots\longrightarrow0
\]
and at sampled times
\[
  t_0<t_1<t_2<\cdots\uparrow T_*<\infty,
  \qquad
  t_{n+1}-t_n\asymp\tau_n,
\]
with comparison constants independent of \(n\). The remaining time would then satisfy
\[
  T_*-t_n
  \asymp
  \sum_{k=n}^{\infty}\tau_k
  =
  \frac{4r_n^2}{3\nu}.
\]
Along this conditional one-field history, the next scale-change interval and the total time left are both comparable to \(r_n^2/\nu\). Their collapse turns the spatial cascade into a temporal demand on the same field, while that field's first and successive Eulerian time derivatives at one fixed spatial point remain uncontrolled.
